\section{Introduction}
\label{sec:introduction}

The transport of hydrogen through metallic and polymeric barriers is of central
importance in energy systems, chemical processing, and materials engineering.
In gas-driven permeation experiments, a membrane of thickness~$L$ separates a
high-pressure feed chamber from a low-pressure permeate chamber.  Hydrogen
dissolves into the upstream face of the membrane according to Sieverts' law
(metals) or Henry's law (polymers), diffuses down the resulting concentration
gradient, and emerges at the downstream face.  The cumulative quantity of gas
that has permeated, $\Cum(t)$, is recorded as a function of time.

\subsection{The Timelag Method}

After an initial transient, $\Cum(t)$ approaches a linear asymptote whose slope
equals the steady-state flux~$\Jss$ and whose $t$-intercept defines the
\emph{timelag}~$\tlag$.  This elegant technique, introduced by
\textcite{Daynes1920} and systematised by \textcite{Barrer1951}, provides a
direct route to the diffusion coefficient.  For a homogeneous membrane with
constant diffusivity~$D$, the classical result is
\begin{equation}
  \label{eq:classical-timelag}
  \tlag = \frac{L^{2}}{6D},
\end{equation}
requiring only the membrane thickness and the measured intercept.

\subsection{Industrial Motivation: Multilayer Barriers}

Many practical barriers are not homogeneous single layers.  Metallic pipelines
may carry oxide scales or deliberately applied coatings; polymer liners in
composite pressure vessels can be laminated with barrier films; and cladding
strategies in nuclear fuel systems create multilayer diffusion paths.  In each
case, hydrogen must traverse two or more distinct material layers, each
characterised by its own diffusivity~$\Ds{i}$ and solubility~$\Ss{i}$.

When a timelag measurement is performed on such a composite membrane, the
resulting $\tlag$ reflects an \emph{effective} combination of the individual
layer properties.  The fundamental question addressed in this work is:

\begin{quote}
\itshape
Given experimental permeation data on a bilayer membrane whose layer
thicknesses are known but whose individual diffusivities and solubilities are
\emph{not}, can the individual transport parameters $\Ds{1}$, $\Ds{2}$,
$\Ss{1}$, $\Ss{2}$ be uniquely extracted?
\end{quote}

\subsection{Prior Work}

The study of diffusion timelags in heterogeneous media has a rich history.
\textcite{Frisch1957} provided early theoretical results on timelags for
membranes with position-dependent diffusivity using variational and moment
methods.  \textcite{AshBarrerPetropoulos1963} derived timelag expressions for
laminated slabs, establishing the framework for multilayer analysis.
\textcite{Petropoulos1973} extended the quantitative analysis to glassy
polymers.  \textcite{Crank1975} provided comprehensive treatments of diffusion
in composite media.  Most recently, \textcite{TranMeldonSontag2021} gave a
thorough treatment of transient diffusion in bilayer composites, including the
important finding that the \emph{downstream} timelag is independent of the
direction of permeation when the boundary conditions are ideal (zero
concentration at the permeate face).

\subsection{Contributions and Outline}

In this paper we provide:
\begin{enumerate}
  \item A self-contained, rigorous derivation of the single-layer timelag via
        Laplace transforms (\cref{sec:single-layer}), establishing the
        methodology.
  \item A complete derivation of the bilayer composite timelag from the
        coupled diffusion equations (\cref{sec:two-layer}), including:
        \begin{itemize}
          \item the full Laplace-domain solution for both layers,
          \item the small-$s$ expansion that yields the timelag,
          \item a proof that the downstream timelag is direction-independent,
          \item derivation of the upstream lead time, which \emph{is}
                direction-dependent.
        \end{itemize}
  \item Systematic extraction methods for the individual transport parameters
        $\Ds{1}$, $\Ds{2}$, $\Ss{1}$, $\Ss{2}$ from experimentally
        accessible quantities (\cref{sec:extraction}), addressing both the
        case where solubilities are known and the general case where they are
        not, with analysis of conditioning and practical protocols.
\end{enumerate}

Detailed algebraic manipulations are collected in \cref{sec:appendix-laplace}
to keep the main exposition readable.
